\section{Future Improvements}

Several improvements can be made to enhance both the functionality and technical depth of the system. 
Future development could focus on expanding the dataset to include a wider range of geographic regions 
and a more diverse selection of cuisines. 
Integrating a more robust and internally managed database would reduce dependency on third-party sources 
and improve long-term reliability, consistency, and scalability.

From a feature perspective, the system could be enhanced by incorporating user authentication 
and personalization mechanisms. 
Allowing users to create accounts, save preferences, and receive customized recommendations 
would significantly improve engagement. 
Additional features such as review moderation, cuisine ranking algorithms, 
and real-time updates could further strengthen the interactivity and usability of the application.
Search functions in particular would be the feature to implement as a priority,
as they would allow users to quickly locate specific regions or cuisines without navigating through multiple pages.
Integrating visual content, including images of dishes, would also improve the overall user experience, 
as food-related platforms benefit greatly from visual representation.

From a technical and project perspective, future work could involve implementing and systematically 
comparing different architectural and database design choices. 
For example, performance metrics such as response time and load time could be measured 
to evaluate a monolithic architecture against a distributed system design. 
Similarly, caching strategies could be compared against a non-cached implementation 
to analyze their impact on latency and scalability under increasing load.

Further experimentation could include evaluating database optimization strategies, 
such as partitioning versus non-partitioned schemas, indexing versus non-indexed queries, 
and the impact of resolving the N+1 query problem compared to leaving it unoptimized. 
Conducting controlled performance benchmarks across these configurations 
would provide deeper empirical insight into backend trade-offs 
and strengthen the analytical contribution of the project.

Overall, these improvements would extend the system beyond functional implementation 
and allow for a more thorough evaluation of backend architectural design choices.